% =============================================================
%  FRONTESPIZIO TEMPLATE  –  stile accademico classico
%  Fedele al layout: Technical Research Report (Duyar 2026)
%  Uso: sostituire le macro nella sezione "DATI DOCUMENTO"
% =============================================================
\documentclass[12pt, a4paper]{article}

\usepackage[T1]{fontenc}
\usepackage[utf8]{inputenc}
\usepackage[english]{babel}
\usepackage{geometry}
\usepackage{amsmath, amssymb}
\usepackage{graphicx}
\usepackage{hyperref}
\usepackage{tabularx}
\usepackage{ragged2e}

\geometry{
    a4paper,
    top    = 3.0cm,
    bottom = 3.0cm,
    left   = 3.5cm,
    right  = 3.5cm
}

\hypersetup{
    colorlinks = true,
    urlcolor   = blue,
    linkcolor  = black,
    citecolor  = black
}

% Sopprime numero pagina sul frontespizio
\pagenumbering{gobble}

% =============================================================
%  DATI DOCUMENTO — modifica solo questa sezione
% =============================================================

\newcommand{\DocCategory}{Technical Research Report}
\newcommand{\DocField}{Power Electronics \& Control Systems}

\newcommand{\DocTitle}{%
    DAB DC/DC Converter for BESS\\
    Applications
}
\newcommand{\DocSubtitle}{%
    Design, Simulation and Experimental Validation of a 250\,kW\\
    Dual-Active-Bridge at 12\,kHz Switching Frequency
}

\newcommand{\DocAuthor}{Davide Bagnara}
\newcommand{\DocDate}{May 2026}

\newcommand{\DocAbstract}{%
    This report presents the design and analysis of a 250\,kW
    Dual-Active-Bridge (DAB) DC/DC converter targeting Battery Energy
    Storage System (BESS) applications, scalable to 1\,MW through parallel
    module stacking. The modulation strategy is investigated through
    Single-Phase-Shift (SPS), Extended-Phase-Shift (EPS), and
    Dual-Phase-Shift (DPS) control laws, with ZVS boundary conditions
    derived analytically for the SKM1700MB20R4S2I4 SiC module.
    Transformer core losses are evaluated via the improved Generalized
    Steinmetz Equation (iGSE). The proposed framework achieves full-load
    efficiency above 98.2\% and guarantees ZVS operation over a dc-link
    voltage range of 600--1000\,V.
}

% Tre righe metadati in fondo (label small caps + valore)
\newcommand{\MetaLabelA}{Project Repository}
\newcommand{\MetaValueA}{https://github.com/pwr-control/dab-bess}

\newcommand{\MetaLabelB}{Simulation Tool}
\newcommand{\MetaValueB}{MATLAB/Simulink R2024b + Simscape Electrical}

\newcommand{\MetaLabelC}{Hardware Platform}
\newcommand{\MetaValueC}{STM32G474 + Semikron SKM1700MB20R4S2I4 SiC}

% =============================================================
%  FRONTESPIZIO
% =============================================================
\begin{document}

% ---- Intestazione centrata ----
\begin{center}
    {\large\scshape \DocCategory}\\[4pt]
    {\small\itshape \DocField}
\end{center}

\vspace{10pt}
\noindent\rule{\linewidth}{0.8pt}

% ---- Titolo ----
\vspace{28pt}
\begin{center}
    {\fontsize{22}{30}\selectfont\bfseries \DocTitle}
\end{center}

\vspace{4pt}
\noindent\rule{\linewidth}{0.8pt}

% ---- Sottotitolo ----
\vspace{14pt}
\begin{center}
    {\large\itshape \DocSubtitle}
\end{center}

% ---- Spazio prima di Autore/Data ----
\vspace{36pt}

% ---- Author / Date su due colonne ----
\noindent
\begin{minipage}[t]{0.48\linewidth}
    \textbf{Author:}\\[2pt]
    \DocAuthor
\end{minipage}%
\hfill
\begin{minipage}[t]{0.48\linewidth}
    \raggedleft
    \textbf{Date:}\\[2pt]
    \DocDate
\end{minipage}

% ---- Spazio prima dell'abstract ----
\vspace{36pt}

% ---- Abstract indentato ----
\begin{center}
\begin{minipage}{0.88\linewidth}
    \small
    \textbf{Abstract:} \DocAbstract
\end{minipage}
\end{center}

% ---- Spazio prima dei metadati ----
\vfill

% ---- Metadati in fondo ----
\noindent
\begin{tabular}{@{} r @{\quad} l @{}}
    \textsc{\small\MetaLabelA:} & {\small\url{\MetaValueA}} \\[3pt]
    \textsc{\small\MetaLabelB:} & {\small\MetaValueB}       \\[3pt]
    \textsc{\small\MetaLabelC:} & {\small\MetaValueC}       \\
\end{tabular}

% =============================================================
%  Fine frontespizio
% =============================================================
\newpage
\pagenumbering{arabic}
\setcounter{page}{1}

% \tableofcontents
% \newpage
% \section{Introduction}

\end{document}
